\section{Joint and Multi-Task Clustering}\label{sec:related:joint}

When two tasks share structure, a common pattern is a shared trunk with task-specific heads, trained by a weighted sum of losses \parencite{caruana1997multitask}.
Image models often attach two softmax heads to one backbone, for example object class and attribute.
That layout assumes both label sets are globally consistent catalogues.
It does not by itself produce clusterable embeddings, and it cannot represent emitter identifiers that are comparable only within a recording.

Hierarchical clustering learns several partitions, but they are nested levels of one relation, not two different relations.
\textcite{yang2016jule} alternate agglomerative merges with representation updates under a triplet loss, so coarser groups are formed from finer ones of the same kind.
A hierarchy of that form would be appropriate if modes were subtypes of emitters, or emitters subtypes of modes.
They are not.
One emitter uses several modes, and one mode appears on several emitters (\cref{sec:method:formulation}).
The two partitions also do not share a label scope: emitter symbols are recording-local, mode symbols are global.

A method that treats the two as softmax heads on a shared backbone therefore asks the wrong questions at inference.
A method that learns a single embedding must compromise between two incompatible similarity relations: pulses that should be close as modes may need to be far as emitters, and conversely.
The dual-branch encoder of \cref{sec:method:arch} is the direct response to that conflict.
The two supervised contrastive terms of \cref{sec:method:loss} differ only in the scope of the positive set, and each branch is decoded by clustering rather than by a closed-catalogue classifier.

\Cref{ch:method} therefore uses a transformer encoder on interleaved \gls{pdw} streams, two contrastive terms with recording-local versus global positives, and independent density clustering of each branch at inference.
